% ============================================================================
% Latent-CoT-Motion: Dataset & Method Sections
% ============================================================================
% Target venue: ECCV (Springer LNCS format)
% Usage: % ============================================================================
% Latent-CoT-Motion: Dataset & Method Sections
% ============================================================================
% Target venue: ECCV (Springer LNCS format)
% Usage: % ============================================================================
% Latent-CoT-Motion: Dataset & Method Sections
% ============================================================================
% Target venue: ECCV (Springer LNCS format)
% Usage: % ============================================================================
% Latent-CoT-Motion: Dataset & Method Sections
% ============================================================================
% Target venue: ECCV (Springer LNCS format)
% Usage: \input{paper_sections.tex} in your main .tex file
%
% Requires in preamble:
%   \usepackage{amsmath,amssymb}
%   \usepackage{booktabs}
% ============================================================================

% ============================================================================
\section{Dataset Construction}
\label{sec:dataset}
% ============================================================================

A key challenge in applying latent reasoning to motion generation lies in
the scarcity of paired data that simultaneously provides textual
descriptions, structured reasoning processes, and ground-truth motion
sequences.
To address this, we construct a large-scale training corpus by integrating
three complementary data sources on top of the
HumanML3D~\cite{guo2022generating} benchmark, as described below.

\subsection{Data Sources and Preprocessing}
\label{sec:data_sources}

\noindent\textbf{Textual motion descriptions.}
We leverage the caption annotations from HumanML3D, which provides
natural-language descriptions paired with 3D motion clips.
Each motion clip may be associated with multiple captions; every caption
line further carries optional temporal boundary tags
$(f_{\mathrm{tag}}, t_{\mathrm{tag}})$ that delineate the relevant
temporal segment, allowing fine-grained caption-to-motion alignment at
the sub-clip level.

\noindent\textbf{Chain-of-thought reasoning annotations.}
To equip the model with structured intermediate reasoning, we adopt the
multi-step thinking annotations from Motion-R1\@.
These annotations decompose the motion generation task into a sequence of
atomic reasoning steps---for instance, \emph{``the person shifts weight
onto the left foot''}---that bridge the semantic gap between abstract
textual descriptions and low-level motion dynamics.
Each sample comprises 2 to 20 steps (mean ${\approx}\,4.2$), which serve
as the explicit chain-of-thought that our curriculum will progressively
compress into latent representations.

\noindent\textbf{Discrete motion tokenisation via VQ-VAE\@.}
Rather than operating directly on continuous joint trajectories, we
discretise motions into a compact token vocabulary amenable to
autoregressive language modelling.
Specifically, we employ a pre-trained VQ-VAE~\cite{zhang2023generating}
with codebook size $K{=}512$ and temporal downsampling factor $d{=}4$
to encode raw 263-dimensional joint-vector sequences into sequences of
discrete codes.
Motions shorter than 40 frames or longer than 200 frames are discarded to
ensure sufficient motion complexity while remaining within the model's
context window, yielding 10--50 motion tokens per sample.

\subsection{Motion Tokenisation}
\label{sec:vqvae}

We follow the VQ-VAE architecture of T2M-GPT~\cite{zhang2023generating},
which employs a 1D convolutional encoder--decoder with an EMA-updated
codebook ($\mu{=}0.99$).
Given a raw motion sequence $\mathbf{m}\in\mathbb{R}^{L\times 263}$
under the HumanML3D 22-joint representation, we first apply z-score
normalisation with precomputed dataset statistics
$(\boldsymbol{\mu},\boldsymbol{\sigma})$ and then quantise:
%
\begin{equation}
  \mathbf{c}
  = \mathrm{VQ\text{-}Enc}\!\bigg(\frac{\mathbf{m}-\boldsymbol{\mu}}
    {\boldsymbol{\sigma}}\bigg),
  \quad c_i\in\{0,1,\ldots,K{-}1\},
  \label{eq:vqvae_enc}
\end{equation}
%
producing a code sequence $\mathbf{c}=(c_1,\ldots,c_T)$ of length
$T{=}\lfloor L/d\rfloor$.
Each code $c_i$ is then mapped to a dedicated token
\texttt{<Motion\_$c_i$>}, and the complete motion is delimited by
\texttt{<Motion>} and \texttt{</Motion>} boundary markers.

\subsection{Training Sample Formulation}
\label{sec:data_format}

Each training instance is assembled into the Qwen ChatML conversational
template and comprises three segments:
(1)~a \textbf{question}~$\mathbf{q}$ consisting of a task-specific system
prompt and the user's motion description, enclosed in
\texttt{<|im\_start|>}/\texttt{<|im\_end|>} delimiters;
(2)~\textbf{reasoning steps}~$\mathbf{s}{=}(s_1,\ldots,s_N)$ wrapped in
\texttt{<think>}$\ldots$\texttt{</think>} tags, which are subject to
progressive latent replacement during curriculum training
(Sec.~\ref{sec:curriculum}); and
(3)~an \textbf{answer}~$\mathbf{a}$ containing the discretised motion
token sequence.
This unified formulation allows us to treat motion generation as a
conditional language modelling problem while retaining explicit reasoning
structure for the curriculum.

\subsection{Dataset Statistics}

Table~\ref{tab:dataset_stats} summarises the resulting dataset, which
follows the standard HumanML3D train/validation/test partition.

\begin{table}[t]
\centering
\caption{Dataset statistics following the HumanML3D partition.}
\label{tab:dataset_stats}
\setlength{\tabcolsep}{8pt}
\begin{tabular}{lccc}
\toprule
 & Train & Val & Test \\
\midrule
Samples & 66\,515 & 4\,170 & 12\,503 \\
Avg.\ motion tokens / sample & \multicolumn{3}{c}{${\sim}37.5$} \\
Avg.\ reasoning steps / sample & \multicolumn{3}{c}{${\sim}4.2$ (range 2--20)} \\
Motion codebook size $K$ & \multicolumn{3}{c}{512} \\
\bottomrule
\end{tabular}
\end{table}


% ============================================================================
\section{Method}
\label{sec:method}
% ============================================================================

\subsection{Overview}

Recent work has demonstrated that large language models can effectively
generate discrete motion token sequences when guided by explicit
chain-of-thought (CoT) reasoning~\cite{zhang2023generating}.
However, producing verbose textual reasoning at inference time introduces
substantial latency and exposes the model to error accumulation over long
reasoning chains.
Inspired by the Coconut framework~\cite{hao2024training}, which showed
that language models can learn to reason in continuous latent spaces
without sacrificing performance, we propose \textbf{Latent-CoT-Motion}---a
training framework that progressively distills explicit CoT reasoning into
compact continuous latent representations for text-to-motion generation.

The key insight is that the intermediate reasoning steps, while beneficial
for learning, need not be produced as discrete tokens at inference time.
Instead, we train the model through a curriculum that gradually replaces
textual reasoning steps with special \emph{latent tokens} whose embeddings
are populated by the model's own hidden states, effectively internalising
the reasoning process.

Our approach comprises three components:
(1)~vocabulary extension and embedding initialisation
    (\S\ref{sec:vocab});
(2)~a multi-pass forward mechanism with hidden-state replacement that
    enables reasoning in continuous space (\S\ref{sec:coconut_forward});
    and
(3)~a stage-based curriculum that orchestrates the transition from
    explicit to latent reasoning (\S\ref{sec:curriculum}).

\subsection{Vocabulary Extension and Embedding Initialisation}
\label{sec:vocab}

Starting from the pre-trained Qwen2.5-3B-Instruct model, we augment its
tokeniser with three groups of task-specific tokens:
(i)~514 \emph{motion tokens} (\texttt{<Motion>}, \texttt{</Motion>}, and
\texttt{<Motion\_0>} through \texttt{<Motion\_511>}) that establish a
one-to-one mapping with the VQ-VAE codebook;
(ii)~2 \emph{think markers} (\texttt{<think>}, \texttt{</think>}) that
delimit explicit reasoning regions; and
(iii)~3 \emph{latent tokens} (\texttt{<|start-latent|>},
\texttt{<|end-latent|>}, \texttt{<|latent|>}) that demarcate the
continuous reasoning region.
The model's embedding matrix and language modelling head are jointly
resized to accommodate the expanded vocabulary.
Following the initialisation strategy of
Coconut~\cite{hao2024training}, we initialise the \texttt{<|latent|>}
embedding from that of ``\texttt{<<}''---a token whose semantic role as a
delimiter provides a more informative starting point than random
initialisation.

\subsection{Multi-Pass Forward with Latent Token Replacement}
\label{sec:coconut_forward}

The central mechanism of our approach replaces the standard single-pass
forward computation with a multi-pass procedure that iteratively fills
latent token embeddings with contextualised hidden states.
Concretely, let
$\mathbf{x}{=}(\mathbf{q},\,t_{\mathrm{s}},\,
  \underbrace{l_1,\ldots,l_M}_{\text{latent}},\,
  t_{\mathrm{e}},\,\mathbf{r},\,\mathbf{a})$
denote an input sequence containing $M$ latent tokens, where $\mathbf{q}$
is the question, $t_{\mathrm{s}}/t_{\mathrm{e}}$ are the latent-region
boundary markers, $\mathbf{r}$ represents any remaining explicit reasoning
steps, and $\mathbf{a}$ is the answer.
The forward computation proceeds as follows.

\paragraph{Iterative hidden-state injection (passes $1,\ldots,M$).}
At the $i$-th pass, the model performs a forward computation from the end
of the previous pass up to the position of the $i$-th latent token $l_i$,
reusing the KV cache accumulated in prior passes to avoid redundant
computation.
The last-layer hidden state $\mathbf{h}_{i-1}$ at the position
immediately preceding $l_i$ is then extracted and used to overwrite the
embedding of $l_i$:
%
\begin{equation}
  \mathbf{e}(l_i) \;\leftarrow\; \mathbf{h}_{i-1}.
  \label{eq:latent_replace}
\end{equation}
%
This operation effectively \emph{conditions} each latent token on the
model's own contextualised representation of all preceding tokens,
allowing information to propagate through the latent region via successive
hidden-state injections rather than through discrete token generation.

\paragraph{Final pass.}
After all $M$ latent embeddings have been replaced, a final forward pass
processes the remaining tokens $(\mathbf{r}, \mathbf{a})$ using the full
accumulated KV cache, producing logits over the vocabulary at each
position.

\paragraph{Loss computation.}
The logits from all $M{+}1$ passes are concatenated along the sequence
dimension:
%
\begin{equation}
  \hat{\mathbf{y}}
  = \big[\hat{\mathbf{y}}^{(1)};\;\hat{\mathbf{y}}^{(2)};\;\ldots\;;\;
    \hat{\mathbf{y}}^{(M+1)}\big],
  \label{eq:logit_concat}
\end{equation}
%
and the model is optimised with the standard next-token cross-entropy
objective:
%
\begin{equation}
  \mathcal{L}
  = -\!\sum_{j\in\mathcal{A}}
    \log\, p_\theta\!\left(x_{j+1}\mid\hat{\mathbf{y}}_j\right),
  \label{eq:loss}
\end{equation}
%
where $\mathcal{A}$ indexes the set of \emph{unmasked} positions.
The question tokens, latent tokens (including boundary markers), and any
padding are assigned label $-100$ and excluded from the loss, so that the
model is supervised exclusively on the remaining reasoning tokens (when
present) and the answer motion tokens.

\paragraph{Efficient batched computation.}
Two design choices ensure computational efficiency.
First, at each pass $i>1$, the KV cache from the preceding pass is sliced
to the appropriate prefix length and reused, avoiding recomputation of
attention over already-processed positions.
Second, we employ a \emph{left-padding collation} strategy that aligns the
first \texttt{<|latent|>} position across all samples in a mini-batch,
enabling the $M$ replacement passes to operate at uniform positions and
thus permitting fully batched multi-pass forward without padding-induced
waste.

\subsection{Curriculum Learning Schedule}
\label{sec:curriculum}

Training the model to reason entirely in latent space from scratch is
intractable, as the latent tokens carry no semantic content at
initialisation.
We therefore adopt a stage-based curriculum that begins with fully explicit
reasoning and gradually increases the proportion of latent tokens.

Formally, let $E_s$ denote the number of training epochs per curriculum
stage and $K_{\max}$ the maximum number of reasoning steps eligible for
replacement.
At epoch $e$, the scheduled curriculum stage is:
%
\begin{equation}
  k(e) = \min\!\Big(\big\lfloor e\,/\,E_s \big\rfloor,\;K_{\max}\Big).
  \label{eq:stage}
\end{equation}
%
At stage $k$, the first $k$ reasoning steps $(s_1,\ldots,s_k)$ are
excised from the input and replaced by $k{\times}c$ latent tokens, where
$c$ denotes the number of latent tokens allocated per replaced step.
The remaining steps $(s_{k+1},\ldots,s_N)$ are retained as explicit text:
%
\begin{equation}
  \mathbf{x}^{(k)}
  = \big(\mathbf{q},\;t_{\mathrm{s}},\;
    \underbrace{l_1,\ldots,l_{kc}}_{\text{latent}},\;
    t_{\mathrm{e}},\;
    \underbrace{s_{k+1}\cdots s_N}_{\text{explicit CoT}},\;
    \mathbf{a}\big).
  \label{eq:curriculum_seq}
\end{equation}
%
Once $k$ reaches $K_{\max}$, all explicit reasoning is removed and the
latent region is zero-padded to a fixed length of $K_{\max}{\cdot}c$
tokens, yielding a fully latent input.
This gradual transition provides the model with a structured learning
signal: in early stages the explicit reasoning tokens act as
``scaffolding'' that guides the latent representations, while in later
stages the model must rely entirely on its internalised continuous
reasoning.

\paragraph{Stage randomisation for regularisation.}
To prevent the model from overfitting to the rigid stage boundaries and to
encourage robustness across varying ratios of latent-to-explicit tokens,
we introduce stochastic stage sampling: with probability $p_u$ each
training sample independently draws a random stage
$k'{\sim}\mathrm{Uniform}(0, N_i)$ in lieu of the scheduled stage $k(e)$,
where $N_i$ is the number of available reasoning steps for sample $i$.
This acts as a form of data augmentation in the curriculum dimension.

Table~\ref{tab:curriculum} summarises the curriculum progression under our
default hyperparameter setting.

\begin{table}[t]
\centering
\caption{Curriculum schedule with $E_s{=}3$, $K_{\max}{=}8$, $c{=}2$.
  The model transitions from fully explicit CoT to fully latent reasoning
  over 30 training epochs.}
\label{tab:curriculum}
\setlength{\tabcolsep}{6pt}
\begin{tabular}{cccp{4.6cm}}
\toprule
Epochs & Stage $k$ & Latent tokens & Supervised target \\
\midrule
0--2   & 0 & 0  & Full CoT + motion tokens \\
3--5   & 1 & 2  & Steps $(s_2,\ldots)$ + motion \\
6--8   & 2 & 4  & Steps $(s_3,\ldots)$ + motion \\
$\vdots$ & $\vdots$ & $\vdots$ & $\vdots$ \\
21--23 & 7 & 14 & Steps $(s_8,\ldots)$ + motion \\
24--29 & ${\geq}8$ & 16 & Motion tokens only \\
\bottomrule
\end{tabular}
\end{table}

\subsection{Inference}
\label{sec:generation}

At test time, the input is constructed with the maximum number of latent
tokens ($K_{\max}{\cdot}c{=}16$) and no explicit reasoning steps,
corresponding to the fully latent regime of the final curriculum stage.
Generation proceeds in two phases:

\noindent\textbf{Phase~1: Latent reasoning.}
The multi-pass forward mechanism (Sec.~\ref{sec:coconut_forward})
processes all $M$ latent tokens, iteratively injecting hidden-state
representations via Eq.~\eqref{eq:latent_replace}.
Upon completion, the model's KV cache encodes a continuous ``chain of
thought'' that captures the reasoning necessary for motion generation,
without producing any intermediate text tokens.

\noindent\textbf{Phase~2: Autoregressive motion decoding.}
Conditioned on the KV cache from Phase~1, the model autoregressively
generates motion tokens via greedy decoding until either an end-of-sequence
token is emitted or the maximum generation budget (256 tokens) is
exhausted.

Compared to explicit CoT inference, this two-phase procedure replaces $T$
sequential text-generation steps with $M$ incremental KV-cached forward
passes, where typically $M{\ll}T$.
We provide a detailed computational analysis of this trade-off in
Sec.~\ref{sec:cost_analysis}.

\subsection{Computational Cost Analysis}
\label{sec:cost_analysis}

We now provide a formal comparison of the inference-time computational
cost between explicit CoT and latent CoT reasoning.
Consider a Transformer decoder with $L$ layers, hidden dimension~$d$,
and attention head dimension~$d_k$ with $n_h$ heads.
Let $|\mathbf{q}|$ denote the question length, $T$ the number of
explicit thinking tokens in a text CoT chain, $M{=}K_{\max}{\cdot}c$ the
number of latent tokens, and $M_a$ the number of motion tokens to
generate.

\paragraph{Unit cost model.}
Under KV-cache-enabled autoregressive decoding, the cost of one
incremental forward step (processing a single new token against a
cache of length~$S$) is dominated by the feed-forward computation for
moderate sequence lengths ($S\ll d^2/d_k$):
%
\begin{equation}
  F_{\mathrm{step}}
  = 2L\bigl(C_{\mathrm{attn}}(S) + C_{\mathrm{ffn}}\bigr)
  \approx 24\,L\,d^2,
  \label{eq:step_cost}
\end{equation}
%
where $C_{\mathrm{attn}}(S){=}4d^2 + 2n_h d_k S$ accounts for
QKV projections and the attention dot product against the cached keys,
and $C_{\mathrm{ffn}}{\approx}8d^2$ for a standard two-layer FFN\@.
A full forward pass over $n$ tokens without KV cache costs approximately
$n\,F_{\mathrm{step}}$, since each position incurs comparable FFN cost
and the quadratic attention term $O(n^2 d_k n_h)$ remains subdominant
for the sequence lengths typical in our setting
($n{\leq}512,\;d{=}3072$).

\paragraph{Explicit CoT inference.}
The model first prefills the question ($|\mathbf{q}|$ tokens), then
autoregressively generates $T$ thinking tokens followed by $M_a$ motion
tokens:
%
\begin{equation}
  C_{\mathrm{explicit}}
  = \underbrace{|\mathbf{q}|\,F_{\mathrm{step}}}_{\text{prefill}}
  \;+\;
  \underbrace{(T + M_a)\,F_{\mathrm{step}}}_{\text{sequential decoding}}.
  \label{eq:cost_explicit}
\end{equation}

\paragraph{Latent CoT inference.}
With KV cache reuse across the $M$ iterative hidden-state replacement
passes, the latent reasoning phase reduces to $M$ incremental forward
steps (each processing one latent position against the accumulated
cache) rather than $M$ full-sequence passes.
The total cost comprises three terms:
%
\begin{equation}
  C_{\mathrm{latent}}
  = \underbrace{|\mathbf{q}|\,F_{\mathrm{step}}}_{\text{prefill}}
  \;+\;
  \underbrace{(M + 1)\,F_{\mathrm{step}}}_{\text{latent passes}}
  \;+\;
  \underbrace{M_a\,F_{\mathrm{step}}}_{\text{motion decoding}},
  \label{eq:cost_latent}
\end{equation}
%
where the $+1$ accounts for the end-of-latent boundary marker.

\paragraph{Reasoning-phase speedup.}
Since the motion decoding cost $M_a\,F_{\mathrm{step}}$ is shared by both
approaches, the inference advantage of latent CoT derives entirely from
the reasoning phase.
The speedup ratio over explicit CoT reasoning is:
%
\begin{equation}
  \rho = \frac{T}{M + 1}.
  \label{eq:speedup}
\end{equation}
%
In our setting, $M{=}K_{\max}{\cdot}c{=}16$ latent tokens replace an
average of $T{\approx}63$ text thinking tokens
(${\sim}4.2$~steps~$\times$~${\sim}15$~tokens/step), yielding
$\rho{\approx}3.7{\times}$.

\paragraph{Discussion.}
The speedup is governed by the \emph{compression ratio} $T/M$: how many
text tokens each latent token effectively replaces.
This compression is possible because latent representations reside in
continuous $\mathbb{R}^d$ space ($d{=}3072$) rather than in a discrete
vocabulary of size $|\mathcal{V}|$, enabling each hidden state to encode
richer intermediate computations than a single text token.
As noted by Hao~et~al.~\cite{hao2024training}, continuous reasoning can
maintain multiple hypotheses simultaneously (akin to breadth-first
search), whereas text CoT is constrained to a single serial reasoning
chain (depth-first search).

However, for tasks where the explicit reasoning chain is short
($T{\lesssim}M$), the multi-pass overhead approaches or exceeds that of
direct text generation, and the latent formulation confers no speed
advantage.
More broadly, the primary motivation for latent CoT in our setting is
not inference speed per se, but rather the capacity of continuous
representations to capture motion-relevant reasoning structure that may
be difficult to express in discrete token form.

\subsection{Implementation Details}
\label{sec:training_details}

We fine-tune Qwen2.5-3B-Instruct in BFloat16 precision using the AdamW
optimiser.
Training is distributed across multiple GPUs via PyTorch
DistributedDataParallel (DDP); we found that FSDP is incompatible with our
multi-pass forward mechanism due to the variable number of forward calls
per sample triggering NCCL synchronisation deadlocks.
To stabilise training under mixed precision, we apply gradient norm
clipping at 1.0\@.
The optimiser state is reinitialised at each curriculum stage transition to
prevent stale momentum from the previous stage's different loss landscape.

Table~\ref{tab:hyperparams} provides the complete hyperparameter
configuration.

\begin{table}[t]
\centering
\caption{Training hyperparameters.}
\label{tab:hyperparams}
\setlength{\tabcolsep}{10pt}
\begin{tabular}{ll}
\toprule
Hyperparameter & Value \\
\midrule
Base model & Qwen2.5-3B-Instruct \\
Precision & BFloat16 \\
Parallelism & DDP \\
Optimiser & AdamW \\
Learning rate & $1\times10^{-4}$ \\
Weight decay & $1\times10^{-2}$ \\
LR warm-up & 200 steps (linear) \\
Gradient clipping & max norm 1.0 \\
Per-GPU batch size & 16 \\
Gradient accumulation & 8 steps \\
Max sequence length & 512 \\
Total epochs & 30 \\
Epochs per stage ($E_s$) & 3 \\
Max latent stage ($K_{\max}$) & 8 \\
Latent tokens per step ($c$) & 2 \\
Stage randomisation ($p_u$) & 0.1 \\
Motion codebook size ($K$) & 512 \\
\bottomrule
\end{tabular}
\end{table}

\paragraph{Validation protocol.}
We monitor two complementary evaluation signals during training.
\emph{Validation loss} is computed every epoch on a held-out subset of
100~samples using the curriculum dataset of the current stage.
\emph{Generation evaluation} is conducted every 3~epochs: the model
generates motion token sequences in the fully latent regime, which are
then compared against ground-truth sequences via two token-level metrics:
(1)~\emph{exact match}---the fraction of samples for which the predicted
and reference motion token sequences are identical; and
(2)~\emph{token accuracy}---the proportion of correctly predicted tokens
in the overlapping prefix, normalised by the ground-truth sequence length.


% ============================================================================
% BibTeX entries (add to your .bib file)
% ============================================================================
% @article{hao2024training,
%   title     = {Training Large Language Models to Reason in a Continuous
%                Latent Space},
%   author    = {Hao, Shibo and Gu, Sainbayar and Ma, Haotian and
%                Hong, Joshua Jiahua and Wang, Zhen and Wang, Daisy Zhe
%                and Hu, Zhiting},
%   journal   = {arXiv preprint arXiv:2412.06769},
%   year      = {2024}
% }
%
% @inproceedings{guo2022generating,
%   title     = {Generating Diverse and Natural {3D} Human Motions from Text},
%   author    = {Guo, Chuan and Zou, Shihao and Zuo, Xinxin and Wang, Sen
%                and Ji, Tao and Li, Xingyu and Cheng, Li},
%   booktitle = {CVPR},
%   year      = {2022}
% }
%
% @inproceedings{zhang2023generating,
%   title     = {{T2M-GPT}: Generating Human Motion from Textual
%                Descriptions with Discrete Representations},
%   author    = {Zhang, Jianrong and Zhang, Yangsong and Cun, Xiaodong
%                and Huang, Shaoli and Zhang, Yong and Zhao, Hongwei
%                and Lu, Hongtao and Shen, Xi},
%   booktitle = {CVPR},
%   year      = {2023}
% }
 in your main .tex file
%
% Requires in preamble:
%   \usepackage{amsmath,amssymb}
%   \usepackage{booktabs}
% ============================================================================

% ============================================================================
\section{Dataset Construction}
\label{sec:dataset}
% ============================================================================

A key challenge in applying latent reasoning to motion generation lies in
the scarcity of paired data that simultaneously provides textual
descriptions, structured reasoning processes, and ground-truth motion
sequences.
To address this, we construct a large-scale training corpus by integrating
three complementary data sources on top of the
HumanML3D~\cite{guo2022generating} benchmark, as described below.

\subsection{Data Sources and Preprocessing}
\label{sec:data_sources}

\noindent\textbf{Textual motion descriptions.}
We leverage the caption annotations from HumanML3D, which provides
natural-language descriptions paired with 3D motion clips.
Each motion clip may be associated with multiple captions; every caption
line further carries optional temporal boundary tags
$(f_{\mathrm{tag}}, t_{\mathrm{tag}})$ that delineate the relevant
temporal segment, allowing fine-grained caption-to-motion alignment at
the sub-clip level.

\noindent\textbf{Chain-of-thought reasoning annotations.}
To equip the model with structured intermediate reasoning, we adopt the
multi-step thinking annotations from Motion-R1\@.
These annotations decompose the motion generation task into a sequence of
atomic reasoning steps---for instance, \emph{``the person shifts weight
onto the left foot''}---that bridge the semantic gap between abstract
textual descriptions and low-level motion dynamics.
Each sample comprises 2 to 20 steps (mean ${\approx}\,4.2$), which serve
as the explicit chain-of-thought that our curriculum will progressively
compress into latent representations.

\noindent\textbf{Discrete motion tokenisation via VQ-VAE\@.}
Rather than operating directly on continuous joint trajectories, we
discretise motions into a compact token vocabulary amenable to
autoregressive language modelling.
Specifically, we employ a pre-trained VQ-VAE~\cite{zhang2023generating}
with codebook size $K{=}512$ and temporal downsampling factor $d{=}4$
to encode raw 263-dimensional joint-vector sequences into sequences of
discrete codes.
Motions shorter than 40 frames or longer than 200 frames are discarded to
ensure sufficient motion complexity while remaining within the model's
context window, yielding 10--50 motion tokens per sample.

\subsection{Motion Tokenisation}
\label{sec:vqvae}

We follow the VQ-VAE architecture of T2M-GPT~\cite{zhang2023generating},
which employs a 1D convolutional encoder--decoder with an EMA-updated
codebook ($\mu{=}0.99$).
Given a raw motion sequence $\mathbf{m}\in\mathbb{R}^{L\times 263}$
under the HumanML3D 22-joint representation, we first apply z-score
normalisation with precomputed dataset statistics
$(\boldsymbol{\mu},\boldsymbol{\sigma})$ and then quantise:
%
\begin{equation}
  \mathbf{c}
  = \mathrm{VQ\text{-}Enc}\!\bigg(\frac{\mathbf{m}-\boldsymbol{\mu}}
    {\boldsymbol{\sigma}}\bigg),
  \quad c_i\in\{0,1,\ldots,K{-}1\},
  \label{eq:vqvae_enc}
\end{equation}
%
producing a code sequence $\mathbf{c}=(c_1,\ldots,c_T)$ of length
$T{=}\lfloor L/d\rfloor$.
Each code $c_i$ is then mapped to a dedicated token
\texttt{<Motion\_$c_i$>}, and the complete motion is delimited by
\texttt{<Motion>} and \texttt{</Motion>} boundary markers.

\subsection{Training Sample Formulation}
\label{sec:data_format}

Each training instance is assembled into the Qwen ChatML conversational
template and comprises three segments:
(1)~a \textbf{question}~$\mathbf{q}$ consisting of a task-specific system
prompt and the user's motion description, enclosed in
\texttt{<|im\_start|>}/\texttt{<|im\_end|>} delimiters;
(2)~\textbf{reasoning steps}~$\mathbf{s}{=}(s_1,\ldots,s_N)$ wrapped in
\texttt{<think>}$\ldots$\texttt{</think>} tags, which are subject to
progressive latent replacement during curriculum training
(Sec.~\ref{sec:curriculum}); and
(3)~an \textbf{answer}~$\mathbf{a}$ containing the discretised motion
token sequence.
This unified formulation allows us to treat motion generation as a
conditional language modelling problem while retaining explicit reasoning
structure for the curriculum.

\subsection{Dataset Statistics}

Table~\ref{tab:dataset_stats} summarises the resulting dataset, which
follows the standard HumanML3D train/validation/test partition.

\begin{table}[t]
\centering
\caption{Dataset statistics following the HumanML3D partition.}
\label{tab:dataset_stats}
\setlength{\tabcolsep}{8pt}
\begin{tabular}{lccc}
\toprule
 & Train & Val & Test \\
\midrule
Samples & 66\,515 & 4\,170 & 12\,503 \\
Avg.\ motion tokens / sample & \multicolumn{3}{c}{${\sim}37.5$} \\
Avg.\ reasoning steps / sample & \multicolumn{3}{c}{${\sim}4.2$ (range 2--20)} \\
Motion codebook size $K$ & \multicolumn{3}{c}{512} \\
\bottomrule
\end{tabular}
\end{table}


% ============================================================================
\section{Method}
\label{sec:method}
% ============================================================================

\subsection{Overview}

Recent work has demonstrated that large language models can effectively
generate discrete motion token sequences when guided by explicit
chain-of-thought (CoT) reasoning~\cite{zhang2023generating}.
However, producing verbose textual reasoning at inference time introduces
substantial latency and exposes the model to error accumulation over long
reasoning chains.
Inspired by the Coconut framework~\cite{hao2024training}, which showed
that language models can learn to reason in continuous latent spaces
without sacrificing performance, we propose \textbf{Latent-CoT-Motion}---a
training framework that progressively distills explicit CoT reasoning into
compact continuous latent representations for text-to-motion generation.

The key insight is that the intermediate reasoning steps, while beneficial
for learning, need not be produced as discrete tokens at inference time.
Instead, we train the model through a curriculum that gradually replaces
textual reasoning steps with special \emph{latent tokens} whose embeddings
are populated by the model's own hidden states, effectively internalising
the reasoning process.

Our approach comprises three components:
(1)~vocabulary extension and embedding initialisation
    (\S\ref{sec:vocab});
(2)~a multi-pass forward mechanism with hidden-state replacement that
    enables reasoning in continuous space (\S\ref{sec:coconut_forward});
    and
(3)~a stage-based curriculum that orchestrates the transition from
    explicit to latent reasoning (\S\ref{sec:curriculum}).

\subsection{Vocabulary Extension and Embedding Initialisation}
\label{sec:vocab}

Starting from the pre-trained Qwen2.5-3B-Instruct model, we augment its
tokeniser with three groups of task-specific tokens:
(i)~514 \emph{motion tokens} (\texttt{<Motion>}, \texttt{</Motion>}, and
\texttt{<Motion\_0>} through \texttt{<Motion\_511>}) that establish a
one-to-one mapping with the VQ-VAE codebook;
(ii)~2 \emph{think markers} (\texttt{<think>}, \texttt{</think>}) that
delimit explicit reasoning regions; and
(iii)~3 \emph{latent tokens} (\texttt{<|start-latent|>},
\texttt{<|end-latent|>}, \texttt{<|latent|>}) that demarcate the
continuous reasoning region.
The model's embedding matrix and language modelling head are jointly
resized to accommodate the expanded vocabulary.
Following the initialisation strategy of
Coconut~\cite{hao2024training}, we initialise the \texttt{<|latent|>}
embedding from that of ``\texttt{<<}''---a token whose semantic role as a
delimiter provides a more informative starting point than random
initialisation.

\subsection{Multi-Pass Forward with Latent Token Replacement}
\label{sec:coconut_forward}

The central mechanism of our approach replaces the standard single-pass
forward computation with a multi-pass procedure that iteratively fills
latent token embeddings with contextualised hidden states.
Concretely, let
$\mathbf{x}{=}(\mathbf{q},\,t_{\mathrm{s}},\,
  \underbrace{l_1,\ldots,l_M}_{\text{latent}},\,
  t_{\mathrm{e}},\,\mathbf{r},\,\mathbf{a})$
denote an input sequence containing $M$ latent tokens, where $\mathbf{q}$
is the question, $t_{\mathrm{s}}/t_{\mathrm{e}}$ are the latent-region
boundary markers, $\mathbf{r}$ represents any remaining explicit reasoning
steps, and $\mathbf{a}$ is the answer.
The forward computation proceeds as follows.

\paragraph{Iterative hidden-state injection (passes $1,\ldots,M$).}
At the $i$-th pass, the model performs a forward computation from the end
of the previous pass up to the position of the $i$-th latent token $l_i$,
reusing the KV cache accumulated in prior passes to avoid redundant
computation.
The last-layer hidden state $\mathbf{h}_{i-1}$ at the position
immediately preceding $l_i$ is then extracted and used to overwrite the
embedding of $l_i$:
%
\begin{equation}
  \mathbf{e}(l_i) \;\leftarrow\; \mathbf{h}_{i-1}.
  \label{eq:latent_replace}
\end{equation}
%
This operation effectively \emph{conditions} each latent token on the
model's own contextualised representation of all preceding tokens,
allowing information to propagate through the latent region via successive
hidden-state injections rather than through discrete token generation.

\paragraph{Final pass.}
After all $M$ latent embeddings have been replaced, a final forward pass
processes the remaining tokens $(\mathbf{r}, \mathbf{a})$ using the full
accumulated KV cache, producing logits over the vocabulary at each
position.

\paragraph{Loss computation.}
The logits from all $M{+}1$ passes are concatenated along the sequence
dimension:
%
\begin{equation}
  \hat{\mathbf{y}}
  = \big[\hat{\mathbf{y}}^{(1)};\;\hat{\mathbf{y}}^{(2)};\;\ldots\;;\;
    \hat{\mathbf{y}}^{(M+1)}\big],
  \label{eq:logit_concat}
\end{equation}
%
and the model is optimised with the standard next-token cross-entropy
objective:
%
\begin{equation}
  \mathcal{L}
  = -\!\sum_{j\in\mathcal{A}}
    \log\, p_\theta\!\left(x_{j+1}\mid\hat{\mathbf{y}}_j\right),
  \label{eq:loss}
\end{equation}
%
where $\mathcal{A}$ indexes the set of \emph{unmasked} positions.
The question tokens, latent tokens (including boundary markers), and any
padding are assigned label $-100$ and excluded from the loss, so that the
model is supervised exclusively on the remaining reasoning tokens (when
present) and the answer motion tokens.

\paragraph{Efficient batched computation.}
Two design choices ensure computational efficiency.
First, at each pass $i>1$, the KV cache from the preceding pass is sliced
to the appropriate prefix length and reused, avoiding recomputation of
attention over already-processed positions.
Second, we employ a \emph{left-padding collation} strategy that aligns the
first \texttt{<|latent|>} position across all samples in a mini-batch,
enabling the $M$ replacement passes to operate at uniform positions and
thus permitting fully batched multi-pass forward without padding-induced
waste.

\subsection{Curriculum Learning Schedule}
\label{sec:curriculum}

Training the model to reason entirely in latent space from scratch is
intractable, as the latent tokens carry no semantic content at
initialisation.
We therefore adopt a stage-based curriculum that begins with fully explicit
reasoning and gradually increases the proportion of latent tokens.

Formally, let $E_s$ denote the number of training epochs per curriculum
stage and $K_{\max}$ the maximum number of reasoning steps eligible for
replacement.
At epoch $e$, the scheduled curriculum stage is:
%
\begin{equation}
  k(e) = \min\!\Big(\big\lfloor e\,/\,E_s \big\rfloor,\;K_{\max}\Big).
  \label{eq:stage}
\end{equation}
%
At stage $k$, the first $k$ reasoning steps $(s_1,\ldots,s_k)$ are
excised from the input and replaced by $k{\times}c$ latent tokens, where
$c$ denotes the number of latent tokens allocated per replaced step.
The remaining steps $(s_{k+1},\ldots,s_N)$ are retained as explicit text:
%
\begin{equation}
  \mathbf{x}^{(k)}
  = \big(\mathbf{q},\;t_{\mathrm{s}},\;
    \underbrace{l_1,\ldots,l_{kc}}_{\text{latent}},\;
    t_{\mathrm{e}},\;
    \underbrace{s_{k+1}\cdots s_N}_{\text{explicit CoT}},\;
    \mathbf{a}\big).
  \label{eq:curriculum_seq}
\end{equation}
%
Once $k$ reaches $K_{\max}$, all explicit reasoning is removed and the
latent region is zero-padded to a fixed length of $K_{\max}{\cdot}c$
tokens, yielding a fully latent input.
This gradual transition provides the model with a structured learning
signal: in early stages the explicit reasoning tokens act as
``scaffolding'' that guides the latent representations, while in later
stages the model must rely entirely on its internalised continuous
reasoning.

\paragraph{Stage randomisation for regularisation.}
To prevent the model from overfitting to the rigid stage boundaries and to
encourage robustness across varying ratios of latent-to-explicit tokens,
we introduce stochastic stage sampling: with probability $p_u$ each
training sample independently draws a random stage
$k'{\sim}\mathrm{Uniform}(0, N_i)$ in lieu of the scheduled stage $k(e)$,
where $N_i$ is the number of available reasoning steps for sample $i$.
This acts as a form of data augmentation in the curriculum dimension.

Table~\ref{tab:curriculum} summarises the curriculum progression under our
default hyperparameter setting.

\begin{table}[t]
\centering
\caption{Curriculum schedule with $E_s{=}3$, $K_{\max}{=}8$, $c{=}2$.
  The model transitions from fully explicit CoT to fully latent reasoning
  over 30 training epochs.}
\label{tab:curriculum}
\setlength{\tabcolsep}{6pt}
\begin{tabular}{cccp{4.6cm}}
\toprule
Epochs & Stage $k$ & Latent tokens & Supervised target \\
\midrule
0--2   & 0 & 0  & Full CoT + motion tokens \\
3--5   & 1 & 2  & Steps $(s_2,\ldots)$ + motion \\
6--8   & 2 & 4  & Steps $(s_3,\ldots)$ + motion \\
$\vdots$ & $\vdots$ & $\vdots$ & $\vdots$ \\
21--23 & 7 & 14 & Steps $(s_8,\ldots)$ + motion \\
24--29 & ${\geq}8$ & 16 & Motion tokens only \\
\bottomrule
\end{tabular}
\end{table}

\subsection{Inference}
\label{sec:generation}

At test time, the input is constructed with the maximum number of latent
tokens ($K_{\max}{\cdot}c{=}16$) and no explicit reasoning steps,
corresponding to the fully latent regime of the final curriculum stage.
Generation proceeds in two phases:

\noindent\textbf{Phase~1: Latent reasoning.}
The multi-pass forward mechanism (Sec.~\ref{sec:coconut_forward})
processes all $M$ latent tokens, iteratively injecting hidden-state
representations via Eq.~\eqref{eq:latent_replace}.
Upon completion, the model's KV cache encodes a continuous ``chain of
thought'' that captures the reasoning necessary for motion generation,
without producing any intermediate text tokens.

\noindent\textbf{Phase~2: Autoregressive motion decoding.}
Conditioned on the KV cache from Phase~1, the model autoregressively
generates motion tokens via greedy decoding until either an end-of-sequence
token is emitted or the maximum generation budget (256 tokens) is
exhausted.

Compared to explicit CoT inference, this two-phase procedure replaces $T$
sequential text-generation steps with $M$ incremental KV-cached forward
passes, where typically $M{\ll}T$.
We provide a detailed computational analysis of this trade-off in
Sec.~\ref{sec:cost_analysis}.

\subsection{Computational Cost Analysis}
\label{sec:cost_analysis}

We now provide a formal comparison of the inference-time computational
cost between explicit CoT and latent CoT reasoning.
Consider a Transformer decoder with $L$ layers, hidden dimension~$d$,
and attention head dimension~$d_k$ with $n_h$ heads.
Let $|\mathbf{q}|$ denote the question length, $T$ the number of
explicit thinking tokens in a text CoT chain, $M{=}K_{\max}{\cdot}c$ the
number of latent tokens, and $M_a$ the number of motion tokens to
generate.

\paragraph{Unit cost model.}
Under KV-cache-enabled autoregressive decoding, the cost of one
incremental forward step (processing a single new token against a
cache of length~$S$) is dominated by the feed-forward computation for
moderate sequence lengths ($S\ll d^2/d_k$):
%
\begin{equation}
  F_{\mathrm{step}}
  = 2L\bigl(C_{\mathrm{attn}}(S) + C_{\mathrm{ffn}}\bigr)
  \approx 24\,L\,d^2,
  \label{eq:step_cost}
\end{equation}
%
where $C_{\mathrm{attn}}(S){=}4d^2 + 2n_h d_k S$ accounts for
QKV projections and the attention dot product against the cached keys,
and $C_{\mathrm{ffn}}{\approx}8d^2$ for a standard two-layer FFN\@.
A full forward pass over $n$ tokens without KV cache costs approximately
$n\,F_{\mathrm{step}}$, since each position incurs comparable FFN cost
and the quadratic attention term $O(n^2 d_k n_h)$ remains subdominant
for the sequence lengths typical in our setting
($n{\leq}512,\;d{=}3072$).

\paragraph{Explicit CoT inference.}
The model first prefills the question ($|\mathbf{q}|$ tokens), then
autoregressively generates $T$ thinking tokens followed by $M_a$ motion
tokens:
%
\begin{equation}
  C_{\mathrm{explicit}}
  = \underbrace{|\mathbf{q}|\,F_{\mathrm{step}}}_{\text{prefill}}
  \;+\;
  \underbrace{(T + M_a)\,F_{\mathrm{step}}}_{\text{sequential decoding}}.
  \label{eq:cost_explicit}
\end{equation}

\paragraph{Latent CoT inference.}
With KV cache reuse across the $M$ iterative hidden-state replacement
passes, the latent reasoning phase reduces to $M$ incremental forward
steps (each processing one latent position against the accumulated
cache) rather than $M$ full-sequence passes.
The total cost comprises three terms:
%
\begin{equation}
  C_{\mathrm{latent}}
  = \underbrace{|\mathbf{q}|\,F_{\mathrm{step}}}_{\text{prefill}}
  \;+\;
  \underbrace{(M + 1)\,F_{\mathrm{step}}}_{\text{latent passes}}
  \;+\;
  \underbrace{M_a\,F_{\mathrm{step}}}_{\text{motion decoding}},
  \label{eq:cost_latent}
\end{equation}
%
where the $+1$ accounts for the end-of-latent boundary marker.

\paragraph{Reasoning-phase speedup.}
Since the motion decoding cost $M_a\,F_{\mathrm{step}}$ is shared by both
approaches, the inference advantage of latent CoT derives entirely from
the reasoning phase.
The speedup ratio over explicit CoT reasoning is:
%
\begin{equation}
  \rho = \frac{T}{M + 1}.
  \label{eq:speedup}
\end{equation}
%
In our setting, $M{=}K_{\max}{\cdot}c{=}16$ latent tokens replace an
average of $T{\approx}63$ text thinking tokens
(${\sim}4.2$~steps~$\times$~${\sim}15$~tokens/step), yielding
$\rho{\approx}3.7{\times}$.

\paragraph{Discussion.}
The speedup is governed by the \emph{compression ratio} $T/M$: how many
text tokens each latent token effectively replaces.
This compression is possible because latent representations reside in
continuous $\mathbb{R}^d$ space ($d{=}3072$) rather than in a discrete
vocabulary of size $|\mathcal{V}|$, enabling each hidden state to encode
richer intermediate computations than a single text token.
As noted by Hao~et~al.~\cite{hao2024training}, continuous reasoning can
maintain multiple hypotheses simultaneously (akin to breadth-first
search), whereas text CoT is constrained to a single serial reasoning
chain (depth-first search).

However, for tasks where the explicit reasoning chain is short
($T{\lesssim}M$), the multi-pass overhead approaches or exceeds that of
direct text generation, and the latent formulation confers no speed
advantage.
More broadly, the primary motivation for latent CoT in our setting is
not inference speed per se, but rather the capacity of continuous
representations to capture motion-relevant reasoning structure that may
be difficult to express in discrete token form.

\subsection{Implementation Details}
\label{sec:training_details}

We fine-tune Qwen2.5-3B-Instruct in BFloat16 precision using the AdamW
optimiser.
Training is distributed across multiple GPUs via PyTorch
DistributedDataParallel (DDP); we found that FSDP is incompatible with our
multi-pass forward mechanism due to the variable number of forward calls
per sample triggering NCCL synchronisation deadlocks.
To stabilise training under mixed precision, we apply gradient norm
clipping at 1.0\@.
The optimiser state is reinitialised at each curriculum stage transition to
prevent stale momentum from the previous stage's different loss landscape.

Table~\ref{tab:hyperparams} provides the complete hyperparameter
configuration.

\begin{table}[t]
\centering
\caption{Training hyperparameters.}
\label{tab:hyperparams}
\setlength{\tabcolsep}{10pt}
\begin{tabular}{ll}
\toprule
Hyperparameter & Value \\
\midrule
Base model & Qwen2.5-3B-Instruct \\
Precision & BFloat16 \\
Parallelism & DDP \\
Optimiser & AdamW \\
Learning rate & $1\times10^{-4}$ \\
Weight decay & $1\times10^{-2}$ \\
LR warm-up & 200 steps (linear) \\
Gradient clipping & max norm 1.0 \\
Per-GPU batch size & 16 \\
Gradient accumulation & 8 steps \\
Max sequence length & 512 \\
Total epochs & 30 \\
Epochs per stage ($E_s$) & 3 \\
Max latent stage ($K_{\max}$) & 8 \\
Latent tokens per step ($c$) & 2 \\
Stage randomisation ($p_u$) & 0.1 \\
Motion codebook size ($K$) & 512 \\
\bottomrule
\end{tabular}
\end{table}

\paragraph{Validation protocol.}
We monitor two complementary evaluation signals during training.
\emph{Validation loss} is computed every epoch on a held-out subset of
100~samples using the curriculum dataset of the current stage.
\emph{Generation evaluation} is conducted every 3~epochs: the model
generates motion token sequences in the fully latent regime, which are
then compared against ground-truth sequences via two token-level metrics:
(1)~\emph{exact match}---the fraction of samples for which the predicted
and reference motion token sequences are identical; and
(2)~\emph{token accuracy}---the proportion of correctly predicted tokens
in the overlapping prefix, normalised by the ground-truth sequence length.


% ============================================================================
% BibTeX entries (add to your .bib file)
% ============================================================================
% @article{hao2024training,
%   title     = {Training Large Language Models to Reason in a Continuous
%                Latent Space},
%   author    = {Hao, Shibo and Gu, Sainbayar and Ma, Haotian and
%                Hong, Joshua Jiahua and Wang, Zhen and Wang, Daisy Zhe
%                and Hu, Zhiting},
%   journal   = {arXiv preprint arXiv:2412.06769},
%   year      = {2024}
% }
%
% @inproceedings{guo2022generating,
%   title     = {Generating Diverse and Natural {3D} Human Motions from Text},
%   author    = {Guo, Chuan and Zou, Shihao and Zuo, Xinxin and Wang, Sen
%                and Ji, Tao and Li, Xingyu and Cheng, Li},
%   booktitle = {CVPR},
%   year      = {2022}
% }
%
% @inproceedings{zhang2023generating,
%   title     = {{T2M-GPT}: Generating Human Motion from Textual
%                Descriptions with Discrete Representations},
%   author    = {Zhang, Jianrong and Zhang, Yangsong and Cun, Xiaodong
%                and Huang, Shaoli and Zhang, Yong and Zhao, Hongwei
%                and Lu, Hongtao and Shen, Xi},
%   booktitle = {CVPR},
%   year      = {2023}
% }
 in your main .tex file
%
% Requires in preamble:
%   \usepackage{amsmath,amssymb}
%   \usepackage{booktabs}
% ============================================================================

% ============================================================================
\section{Dataset Construction}
\label{sec:dataset}
% ============================================================================

A key challenge in applying latent reasoning to motion generation lies in
the scarcity of paired data that simultaneously provides textual
descriptions, structured reasoning processes, and ground-truth motion
sequences.
To address this, we construct a large-scale training corpus by integrating
three complementary data sources on top of the
HumanML3D~\cite{guo2022generating} benchmark, as described below.

\subsection{Data Sources and Preprocessing}
\label{sec:data_sources}

\noindent\textbf{Textual motion descriptions.}
We leverage the caption annotations from HumanML3D, which provides
natural-language descriptions paired with 3D motion clips.
Each motion clip may be associated with multiple captions; every caption
line further carries optional temporal boundary tags
$(f_{\mathrm{tag}}, t_{\mathrm{tag}})$ that delineate the relevant
temporal segment, allowing fine-grained caption-to-motion alignment at
the sub-clip level.

\noindent\textbf{Chain-of-thought reasoning annotations.}
To equip the model with structured intermediate reasoning, we adopt the
multi-step thinking annotations from Motion-R1\@.
These annotations decompose the motion generation task into a sequence of
atomic reasoning steps---for instance, \emph{``the person shifts weight
onto the left foot''}---that bridge the semantic gap between abstract
textual descriptions and low-level motion dynamics.
Each sample comprises 2 to 20 steps (mean ${\approx}\,4.2$), which serve
as the explicit chain-of-thought that our curriculum will progressively
compress into latent representations.

\noindent\textbf{Discrete motion tokenisation via VQ-VAE\@.}
Rather than operating directly on continuous joint trajectories, we
discretise motions into a compact token vocabulary amenable to
autoregressive language modelling.
Specifically, we employ a pre-trained VQ-VAE~\cite{zhang2023generating}
with codebook size $K{=}512$ and temporal downsampling factor $d{=}4$
to encode raw 263-dimensional joint-vector sequences into sequences of
discrete codes.
Motions shorter than 40 frames or longer than 200 frames are discarded to
ensure sufficient motion complexity while remaining within the model's
context window, yielding 10--50 motion tokens per sample.

\subsection{Motion Tokenisation}
\label{sec:vqvae}

We follow the VQ-VAE architecture of T2M-GPT~\cite{zhang2023generating},
which employs a 1D convolutional encoder--decoder with an EMA-updated
codebook ($\mu{=}0.99$).
Given a raw motion sequence $\mathbf{m}\in\mathbb{R}^{L\times 263}$
under the HumanML3D 22-joint representation, we first apply z-score
normalisation with precomputed dataset statistics
$(\boldsymbol{\mu},\boldsymbol{\sigma})$ and then quantise:
%
\begin{equation}
  \mathbf{c}
  = \mathrm{VQ\text{-}Enc}\!\bigg(\frac{\mathbf{m}-\boldsymbol{\mu}}
    {\boldsymbol{\sigma}}\bigg),
  \quad c_i\in\{0,1,\ldots,K{-}1\},
  \label{eq:vqvae_enc}
\end{equation}
%
producing a code sequence $\mathbf{c}=(c_1,\ldots,c_T)$ of length
$T{=}\lfloor L/d\rfloor$.
Each code $c_i$ is then mapped to a dedicated token
\texttt{<Motion\_$c_i$>}, and the complete motion is delimited by
\texttt{<Motion>} and \texttt{</Motion>} boundary markers.

\subsection{Training Sample Formulation}
\label{sec:data_format}

Each training instance is assembled into the Qwen ChatML conversational
template and comprises three segments:
(1)~a \textbf{question}~$\mathbf{q}$ consisting of a task-specific system
prompt and the user's motion description, enclosed in
\texttt{<|im\_start|>}/\texttt{<|im\_end|>} delimiters;
(2)~\textbf{reasoning steps}~$\mathbf{s}{=}(s_1,\ldots,s_N)$ wrapped in
\texttt{<think>}$\ldots$\texttt{</think>} tags, which are subject to
progressive latent replacement during curriculum training
(Sec.~\ref{sec:curriculum}); and
(3)~an \textbf{answer}~$\mathbf{a}$ containing the discretised motion
token sequence.
This unified formulation allows us to treat motion generation as a
conditional language modelling problem while retaining explicit reasoning
structure for the curriculum.

\subsection{Dataset Statistics}

Table~\ref{tab:dataset_stats} summarises the resulting dataset, which
follows the standard HumanML3D train/validation/test partition.

\begin{table}[t]
\centering
\caption{Dataset statistics following the HumanML3D partition.}
\label{tab:dataset_stats}
\setlength{\tabcolsep}{8pt}
\begin{tabular}{lccc}
\toprule
 & Train & Val & Test \\
\midrule
Samples & 66\,515 & 4\,170 & 12\,503 \\
Avg.\ motion tokens / sample & \multicolumn{3}{c}{${\sim}37.5$} \\
Avg.\ reasoning steps / sample & \multicolumn{3}{c}{${\sim}4.2$ (range 2--20)} \\
Motion codebook size $K$ & \multicolumn{3}{c}{512} \\
\bottomrule
\end{tabular}
\end{table}


% ============================================================================
\section{Method}
\label{sec:method}
% ============================================================================

\subsection{Overview}

Recent work has demonstrated that large language models can effectively
generate discrete motion token sequences when guided by explicit
chain-of-thought (CoT) reasoning~\cite{zhang2023generating}.
However, producing verbose textual reasoning at inference time introduces
substantial latency and exposes the model to error accumulation over long
reasoning chains.
Inspired by the Coconut framework~\cite{hao2024training}, which showed
that language models can learn to reason in continuous latent spaces
without sacrificing performance, we propose \textbf{Latent-CoT-Motion}---a
training framework that progressively distills explicit CoT reasoning into
compact continuous latent representations for text-to-motion generation.

The key insight is that the intermediate reasoning steps, while beneficial
for learning, need not be produced as discrete tokens at inference time.
Instead, we train the model through a curriculum that gradually replaces
textual reasoning steps with special \emph{latent tokens} whose embeddings
are populated by the model's own hidden states, effectively internalising
the reasoning process.

Our approach comprises three components:
(1)~vocabulary extension and embedding initialisation
    (\S\ref{sec:vocab});
(2)~a multi-pass forward mechanism with hidden-state replacement that
    enables reasoning in continuous space (\S\ref{sec:coconut_forward});
    and
(3)~a stage-based curriculum that orchestrates the transition from
    explicit to latent reasoning (\S\ref{sec:curriculum}).

\subsection{Vocabulary Extension and Embedding Initialisation}
\label{sec:vocab}

Starting from the pre-trained Qwen2.5-3B-Instruct model, we augment its
tokeniser with three groups of task-specific tokens:
(i)~514 \emph{motion tokens} (\texttt{<Motion>}, \texttt{</Motion>}, and
\texttt{<Motion\_0>} through \texttt{<Motion\_511>}) that establish a
one-to-one mapping with the VQ-VAE codebook;
(ii)~2 \emph{think markers} (\texttt{<think>}, \texttt{</think>}) that
delimit explicit reasoning regions; and
(iii)~3 \emph{latent tokens} (\texttt{<|start-latent|>},
\texttt{<|end-latent|>}, \texttt{<|latent|>}) that demarcate the
continuous reasoning region.
The model's embedding matrix and language modelling head are jointly
resized to accommodate the expanded vocabulary.
Following the initialisation strategy of
Coconut~\cite{hao2024training}, we initialise the \texttt{<|latent|>}
embedding from that of ``\texttt{<<}''---a token whose semantic role as a
delimiter provides a more informative starting point than random
initialisation.

\subsection{Multi-Pass Forward with Latent Token Replacement}
\label{sec:coconut_forward}

The central mechanism of our approach replaces the standard single-pass
forward computation with a multi-pass procedure that iteratively fills
latent token embeddings with contextualised hidden states.
Concretely, let
$\mathbf{x}{=}(\mathbf{q},\,t_{\mathrm{s}},\,
  \underbrace{l_1,\ldots,l_M}_{\text{latent}},\,
  t_{\mathrm{e}},\,\mathbf{r},\,\mathbf{a})$
denote an input sequence containing $M$ latent tokens, where $\mathbf{q}$
is the question, $t_{\mathrm{s}}/t_{\mathrm{e}}$ are the latent-region
boundary markers, $\mathbf{r}$ represents any remaining explicit reasoning
steps, and $\mathbf{a}$ is the answer.
The forward computation proceeds as follows.

\paragraph{Iterative hidden-state injection (passes $1,\ldots,M$).}
At the $i$-th pass, the model performs a forward computation from the end
of the previous pass up to the position of the $i$-th latent token $l_i$,
reusing the KV cache accumulated in prior passes to avoid redundant
computation.
The last-layer hidden state $\mathbf{h}_{i-1}$ at the position
immediately preceding $l_i$ is then extracted and used to overwrite the
embedding of $l_i$:
%
\begin{equation}
  \mathbf{e}(l_i) \;\leftarrow\; \mathbf{h}_{i-1}.
  \label{eq:latent_replace}
\end{equation}
%
This operation effectively \emph{conditions} each latent token on the
model's own contextualised representation of all preceding tokens,
allowing information to propagate through the latent region via successive
hidden-state injections rather than through discrete token generation.

\paragraph{Final pass.}
After all $M$ latent embeddings have been replaced, a final forward pass
processes the remaining tokens $(\mathbf{r}, \mathbf{a})$ using the full
accumulated KV cache, producing logits over the vocabulary at each
position.

\paragraph{Loss computation.}
The logits from all $M{+}1$ passes are concatenated along the sequence
dimension:
%
\begin{equation}
  \hat{\mathbf{y}}
  = \big[\hat{\mathbf{y}}^{(1)};\;\hat{\mathbf{y}}^{(2)};\;\ldots\;;\;
    \hat{\mathbf{y}}^{(M+1)}\big],
  \label{eq:logit_concat}
\end{equation}
%
and the model is optimised with the standard next-token cross-entropy
objective:
%
\begin{equation}
  \mathcal{L}
  = -\!\sum_{j\in\mathcal{A}}
    \log\, p_\theta\!\left(x_{j+1}\mid\hat{\mathbf{y}}_j\right),
  \label{eq:loss}
\end{equation}
%
where $\mathcal{A}$ indexes the set of \emph{unmasked} positions.
The question tokens, latent tokens (including boundary markers), and any
padding are assigned label $-100$ and excluded from the loss, so that the
model is supervised exclusively on the remaining reasoning tokens (when
present) and the answer motion tokens.

\paragraph{Efficient batched computation.}
Two design choices ensure computational efficiency.
First, at each pass $i>1$, the KV cache from the preceding pass is sliced
to the appropriate prefix length and reused, avoiding recomputation of
attention over already-processed positions.
Second, we employ a \emph{left-padding collation} strategy that aligns the
first \texttt{<|latent|>} position across all samples in a mini-batch,
enabling the $M$ replacement passes to operate at uniform positions and
thus permitting fully batched multi-pass forward without padding-induced
waste.

\subsection{Curriculum Learning Schedule}
\label{sec:curriculum}

Training the model to reason entirely in latent space from scratch is
intractable, as the latent tokens carry no semantic content at
initialisation.
We therefore adopt a stage-based curriculum that begins with fully explicit
reasoning and gradually increases the proportion of latent tokens.

Formally, let $E_s$ denote the number of training epochs per curriculum
stage and $K_{\max}$ the maximum number of reasoning steps eligible for
replacement.
At epoch $e$, the scheduled curriculum stage is:
%
\begin{equation}
  k(e) = \min\!\Big(\big\lfloor e\,/\,E_s \big\rfloor,\;K_{\max}\Big).
  \label{eq:stage}
\end{equation}
%
At stage $k$, the first $k$ reasoning steps $(s_1,\ldots,s_k)$ are
excised from the input and replaced by $k{\times}c$ latent tokens, where
$c$ denotes the number of latent tokens allocated per replaced step.
The remaining steps $(s_{k+1},\ldots,s_N)$ are retained as explicit text:
%
\begin{equation}
  \mathbf{x}^{(k)}
  = \big(\mathbf{q},\;t_{\mathrm{s}},\;
    \underbrace{l_1,\ldots,l_{kc}}_{\text{latent}},\;
    t_{\mathrm{e}},\;
    \underbrace{s_{k+1}\cdots s_N}_{\text{explicit CoT}},\;
    \mathbf{a}\big).
  \label{eq:curriculum_seq}
\end{equation}
%
Once $k$ reaches $K_{\max}$, all explicit reasoning is removed and the
latent region is zero-padded to a fixed length of $K_{\max}{\cdot}c$
tokens, yielding a fully latent input.
This gradual transition provides the model with a structured learning
signal: in early stages the explicit reasoning tokens act as
``scaffolding'' that guides the latent representations, while in later
stages the model must rely entirely on its internalised continuous
reasoning.

\paragraph{Stage randomisation for regularisation.}
To prevent the model from overfitting to the rigid stage boundaries and to
encourage robustness across varying ratios of latent-to-explicit tokens,
we introduce stochastic stage sampling: with probability $p_u$ each
training sample independently draws a random stage
$k'{\sim}\mathrm{Uniform}(0, N_i)$ in lieu of the scheduled stage $k(e)$,
where $N_i$ is the number of available reasoning steps for sample $i$.
This acts as a form of data augmentation in the curriculum dimension.

Table~\ref{tab:curriculum} summarises the curriculum progression under our
default hyperparameter setting.

\begin{table}[t]
\centering
\caption{Curriculum schedule with $E_s{=}3$, $K_{\max}{=}8$, $c{=}2$.
  The model transitions from fully explicit CoT to fully latent reasoning
  over 30 training epochs.}
\label{tab:curriculum}
\setlength{\tabcolsep}{6pt}
\begin{tabular}{cccp{4.6cm}}
\toprule
Epochs & Stage $k$ & Latent tokens & Supervised target \\
\midrule
0--2   & 0 & 0  & Full CoT + motion tokens \\
3--5   & 1 & 2  & Steps $(s_2,\ldots)$ + motion \\
6--8   & 2 & 4  & Steps $(s_3,\ldots)$ + motion \\
$\vdots$ & $\vdots$ & $\vdots$ & $\vdots$ \\
21--23 & 7 & 14 & Steps $(s_8,\ldots)$ + motion \\
24--29 & ${\geq}8$ & 16 & Motion tokens only \\
\bottomrule
\end{tabular}
\end{table}

\subsection{Inference}
\label{sec:generation}

At test time, the input is constructed with the maximum number of latent
tokens ($K_{\max}{\cdot}c{=}16$) and no explicit reasoning steps,
corresponding to the fully latent regime of the final curriculum stage.
Generation proceeds in two phases:

\noindent\textbf{Phase~1: Latent reasoning.}
The multi-pass forward mechanism (Sec.~\ref{sec:coconut_forward})
processes all $M$ latent tokens, iteratively injecting hidden-state
representations via Eq.~\eqref{eq:latent_replace}.
Upon completion, the model's KV cache encodes a continuous ``chain of
thought'' that captures the reasoning necessary for motion generation,
without producing any intermediate text tokens.

\noindent\textbf{Phase~2: Autoregressive motion decoding.}
Conditioned on the KV cache from Phase~1, the model autoregressively
generates motion tokens via greedy decoding until either an end-of-sequence
token is emitted or the maximum generation budget (256 tokens) is
exhausted.

Compared to explicit CoT inference, this two-phase procedure replaces $T$
sequential text-generation steps with $M$ incremental KV-cached forward
passes, where typically $M{\ll}T$.
We provide a detailed computational analysis of this trade-off in
Sec.~\ref{sec:cost_analysis}.

\subsection{Computational Cost Analysis}
\label{sec:cost_analysis}

We now provide a formal comparison of the inference-time computational
cost between explicit CoT and latent CoT reasoning.
Consider a Transformer decoder with $L$ layers, hidden dimension~$d$,
and attention head dimension~$d_k$ with $n_h$ heads.
Let $|\mathbf{q}|$ denote the question length, $T$ the number of
explicit thinking tokens in a text CoT chain, $M{=}K_{\max}{\cdot}c$ the
number of latent tokens, and $M_a$ the number of motion tokens to
generate.

\paragraph{Unit cost model.}
Under KV-cache-enabled autoregressive decoding, the cost of one
incremental forward step (processing a single new token against a
cache of length~$S$) is dominated by the feed-forward computation for
moderate sequence lengths ($S\ll d^2/d_k$):
%
\begin{equation}
  F_{\mathrm{step}}
  = 2L\bigl(C_{\mathrm{attn}}(S) + C_{\mathrm{ffn}}\bigr)
  \approx 24\,L\,d^2,
  \label{eq:step_cost}
\end{equation}
%
where $C_{\mathrm{attn}}(S){=}4d^2 + 2n_h d_k S$ accounts for
QKV projections and the attention dot product against the cached keys,
and $C_{\mathrm{ffn}}{\approx}8d^2$ for a standard two-layer FFN\@.
A full forward pass over $n$ tokens without KV cache costs approximately
$n\,F_{\mathrm{step}}$, since each position incurs comparable FFN cost
and the quadratic attention term $O(n^2 d_k n_h)$ remains subdominant
for the sequence lengths typical in our setting
($n{\leq}512,\;d{=}3072$).

\paragraph{Explicit CoT inference.}
The model first prefills the question ($|\mathbf{q}|$ tokens), then
autoregressively generates $T$ thinking tokens followed by $M_a$ motion
tokens:
%
\begin{equation}
  C_{\mathrm{explicit}}
  = \underbrace{|\mathbf{q}|\,F_{\mathrm{step}}}_{\text{prefill}}
  \;+\;
  \underbrace{(T + M_a)\,F_{\mathrm{step}}}_{\text{sequential decoding}}.
  \label{eq:cost_explicit}
\end{equation}

\paragraph{Latent CoT inference.}
With KV cache reuse across the $M$ iterative hidden-state replacement
passes, the latent reasoning phase reduces to $M$ incremental forward
steps (each processing one latent position against the accumulated
cache) rather than $M$ full-sequence passes.
The total cost comprises three terms:
%
\begin{equation}
  C_{\mathrm{latent}}
  = \underbrace{|\mathbf{q}|\,F_{\mathrm{step}}}_{\text{prefill}}
  \;+\;
  \underbrace{(M + 1)\,F_{\mathrm{step}}}_{\text{latent passes}}
  \;+\;
  \underbrace{M_a\,F_{\mathrm{step}}}_{\text{motion decoding}},
  \label{eq:cost_latent}
\end{equation}
%
where the $+1$ accounts for the end-of-latent boundary marker.

\paragraph{Reasoning-phase speedup.}
Since the motion decoding cost $M_a\,F_{\mathrm{step}}$ is shared by both
approaches, the inference advantage of latent CoT derives entirely from
the reasoning phase.
The speedup ratio over explicit CoT reasoning is:
%
\begin{equation}
  \rho = \frac{T}{M + 1}.
  \label{eq:speedup}
\end{equation}
%
In our setting, $M{=}K_{\max}{\cdot}c{=}16$ latent tokens replace an
average of $T{\approx}63$ text thinking tokens
(${\sim}4.2$~steps~$\times$~${\sim}15$~tokens/step), yielding
$\rho{\approx}3.7{\times}$.

\paragraph{Discussion.}
The speedup is governed by the \emph{compression ratio} $T/M$: how many
text tokens each latent token effectively replaces.
This compression is possible because latent representations reside in
continuous $\mathbb{R}^d$ space ($d{=}3072$) rather than in a discrete
vocabulary of size $|\mathcal{V}|$, enabling each hidden state to encode
richer intermediate computations than a single text token.
As noted by Hao~et~al.~\cite{hao2024training}, continuous reasoning can
maintain multiple hypotheses simultaneously (akin to breadth-first
search), whereas text CoT is constrained to a single serial reasoning
chain (depth-first search).

However, for tasks where the explicit reasoning chain is short
($T{\lesssim}M$), the multi-pass overhead approaches or exceeds that of
direct text generation, and the latent formulation confers no speed
advantage.
More broadly, the primary motivation for latent CoT in our setting is
not inference speed per se, but rather the capacity of continuous
representations to capture motion-relevant reasoning structure that may
be difficult to express in discrete token form.

\subsection{Implementation Details}
\label{sec:training_details}

We fine-tune Qwen2.5-3B-Instruct in BFloat16 precision using the AdamW
optimiser.
Training is distributed across multiple GPUs via PyTorch
DistributedDataParallel (DDP); we found that FSDP is incompatible with our
multi-pass forward mechanism due to the variable number of forward calls
per sample triggering NCCL synchronisation deadlocks.
To stabilise training under mixed precision, we apply gradient norm
clipping at 1.0\@.
The optimiser state is reinitialised at each curriculum stage transition to
prevent stale momentum from the previous stage's different loss landscape.

Table~\ref{tab:hyperparams} provides the complete hyperparameter
configuration.

\begin{table}[t]
\centering
\caption{Training hyperparameters.}
\label{tab:hyperparams}
\setlength{\tabcolsep}{10pt}
\begin{tabular}{ll}
\toprule
Hyperparameter & Value \\
\midrule
Base model & Qwen2.5-3B-Instruct \\
Precision & BFloat16 \\
Parallelism & DDP \\
Optimiser & AdamW \\
Learning rate & $1\times10^{-4}$ \\
Weight decay & $1\times10^{-2}$ \\
LR warm-up & 200 steps (linear) \\
Gradient clipping & max norm 1.0 \\
Per-GPU batch size & 16 \\
Gradient accumulation & 8 steps \\
Max sequence length & 512 \\
Total epochs & 30 \\
Epochs per stage ($E_s$) & 3 \\
Max latent stage ($K_{\max}$) & 8 \\
Latent tokens per step ($c$) & 2 \\
Stage randomisation ($p_u$) & 0.1 \\
Motion codebook size ($K$) & 512 \\
\bottomrule
\end{tabular}
\end{table}

\paragraph{Validation protocol.}
We monitor two complementary evaluation signals during training.
\emph{Validation loss} is computed every epoch on a held-out subset of
100~samples using the curriculum dataset of the current stage.
\emph{Generation evaluation} is conducted every 3~epochs: the model
generates motion token sequences in the fully latent regime, which are
then compared against ground-truth sequences via two token-level metrics:
(1)~\emph{exact match}---the fraction of samples for which the predicted
and reference motion token sequences are identical; and
(2)~\emph{token accuracy}---the proportion of correctly predicted tokens
in the overlapping prefix, normalised by the ground-truth sequence length.


% ============================================================================
% BibTeX entries (add to your .bib file)
% ============================================================================
% @article{hao2024training,
%   title     = {Training Large Language Models to Reason in a Continuous
%                Latent Space},
%   author    = {Hao, Shibo and Gu, Sainbayar and Ma, Haotian and
%                Hong, Joshua Jiahua and Wang, Zhen and Wang, Daisy Zhe
%                and Hu, Zhiting},
%   journal   = {arXiv preprint arXiv:2412.06769},
%   year      = {2024}
% }
%
% @inproceedings{guo2022generating,
%   title     = {Generating Diverse and Natural {3D} Human Motions from Text},
%   author    = {Guo, Chuan and Zou, Shihao and Zuo, Xinxin and Wang, Sen
%                and Ji, Tao and Li, Xingyu and Cheng, Li},
%   booktitle = {CVPR},
%   year      = {2022}
% }
%
% @inproceedings{zhang2023generating,
%   title     = {{T2M-GPT}: Generating Human Motion from Textual
%                Descriptions with Discrete Representations},
%   author    = {Zhang, Jianrong and Zhang, Yangsong and Cun, Xiaodong
%                and Huang, Shaoli and Zhang, Yong and Zhao, Hongwei
%                and Lu, Hongtao and Shen, Xi},
%   booktitle = {CVPR},
%   year      = {2023}
% }
 in your main .tex file
%
% Requires in preamble:
%   \usepackage{amsmath,amssymb}
%   \usepackage{booktabs}
% ============================================================================

% ============================================================================
\section{Dataset Construction}
\label{sec:dataset}
% ============================================================================

A key challenge in applying latent reasoning to motion generation lies in
the scarcity of paired data that simultaneously provides textual
descriptions, structured reasoning processes, and ground-truth motion
sequences.
To address this, we construct a large-scale training corpus by integrating
three complementary data sources on top of the
HumanML3D~\cite{guo2022generating} benchmark, as described below.

\subsection{Data Sources and Preprocessing}
\label{sec:data_sources}

\noindent\textbf{Textual motion descriptions.}
We leverage the caption annotations from HumanML3D, which provides
natural-language descriptions paired with 3D motion clips.
Each motion clip may be associated with multiple captions; every caption
line further carries optional temporal boundary tags
$(f_{\mathrm{tag}}, t_{\mathrm{tag}})$ that delineate the relevant
temporal segment, allowing fine-grained caption-to-motion alignment at
the sub-clip level.

\noindent\textbf{Chain-of-thought reasoning annotations.}
To equip the model with structured intermediate reasoning, we adopt the
multi-step thinking annotations from Motion-R1\@.
These annotations decompose the motion generation task into a sequence of
atomic reasoning steps---for instance, \emph{``the person shifts weight
onto the left foot''}---that bridge the semantic gap between abstract
textual descriptions and low-level motion dynamics.
Each sample comprises 2 to 20 steps (mean ${\approx}\,4.2$), which serve
as the explicit chain-of-thought that our curriculum will progressively
compress into latent representations.

\noindent\textbf{Discrete motion tokenisation via VQ-VAE\@.}
Rather than operating directly on continuous joint trajectories, we
discretise motions into a compact token vocabulary amenable to
autoregressive language modelling.
Specifically, we employ a pre-trained VQ-VAE~\cite{zhang2023generating}
with codebook size $K{=}512$ and temporal downsampling factor $d{=}4$
to encode raw 263-dimensional joint-vector sequences into sequences of
discrete codes.
Motions shorter than 40 frames or longer than 200 frames are discarded to
ensure sufficient motion complexity while remaining within the model's
context window, yielding 10--50 motion tokens per sample.

\subsection{Motion Tokenisation}
\label{sec:vqvae}

We follow the VQ-VAE architecture of T2M-GPT~\cite{zhang2023generating},
which employs a 1D convolutional encoder--decoder with an EMA-updated
codebook ($\mu{=}0.99$).
Given a raw motion sequence $\mathbf{m}\in\mathbb{R}^{L\times 263}$
under the HumanML3D 22-joint representation, we first apply z-score
normalisation with precomputed dataset statistics
$(\boldsymbol{\mu},\boldsymbol{\sigma})$ and then quantise:
%
\begin{equation}
  \mathbf{c}
  = \mathrm{VQ\text{-}Enc}\!\bigg(\frac{\mathbf{m}-\boldsymbol{\mu}}
    {\boldsymbol{\sigma}}\bigg),
  \quad c_i\in\{0,1,\ldots,K{-}1\},
  \label{eq:vqvae_enc}
\end{equation}
%
producing a code sequence $\mathbf{c}=(c_1,\ldots,c_T)$ of length
$T{=}\lfloor L/d\rfloor$.
Each code $c_i$ is then mapped to a dedicated token
\texttt{<Motion\_$c_i$>}, and the complete motion is delimited by
\texttt{<Motion>} and \texttt{</Motion>} boundary markers.

\subsection{Training Sample Formulation}
\label{sec:data_format}

Each training instance is assembled into the Qwen ChatML conversational
template and comprises three segments:
(1)~a \textbf{question}~$\mathbf{q}$ consisting of a task-specific system
prompt and the user's motion description, enclosed in
\texttt{<|im\_start|>}/\texttt{<|im\_end|>} delimiters;
(2)~\textbf{reasoning steps}~$\mathbf{s}{=}(s_1,\ldots,s_N)$ wrapped in
\texttt{<think>}$\ldots$\texttt{</think>} tags, which are subject to
progressive latent replacement during curriculum training
(Sec.~\ref{sec:curriculum}); and
(3)~an \textbf{answer}~$\mathbf{a}$ containing the discretised motion
token sequence.
This unified formulation allows us to treat motion generation as a
conditional language modelling problem while retaining explicit reasoning
structure for the curriculum.

\subsection{Dataset Statistics}

Table~\ref{tab:dataset_stats} summarises the resulting dataset, which
follows the standard HumanML3D train/validation/test partition.

\begin{table}[t]
\centering
\caption{Dataset statistics following the HumanML3D partition.}
\label{tab:dataset_stats}
\setlength{\tabcolsep}{8pt}
\begin{tabular}{lccc}
\toprule
 & Train & Val & Test \\
\midrule
Samples & 66\,515 & 4\,170 & 12\,503 \\
Avg.\ motion tokens / sample & \multicolumn{3}{c}{${\sim}37.5$} \\
Avg.\ reasoning steps / sample & \multicolumn{3}{c}{${\sim}4.2$ (range 2--20)} \\
Motion codebook size $K$ & \multicolumn{3}{c}{512} \\
\bottomrule
\end{tabular}
\end{table}


% ============================================================================
\section{Method}
\label{sec:method}
% ============================================================================

\subsection{Overview}

Recent work has demonstrated that large language models can effectively
generate discrete motion token sequences when guided by explicit
chain-of-thought (CoT) reasoning~\cite{zhang2023generating}.
However, producing verbose textual reasoning at inference time introduces
substantial latency and exposes the model to error accumulation over long
reasoning chains.
Inspired by the Coconut framework~\cite{hao2024training}, which showed
that language models can learn to reason in continuous latent spaces
without sacrificing performance, we propose \textbf{Latent-CoT-Motion}---a
training framework that progressively distills explicit CoT reasoning into
compact continuous latent representations for text-to-motion generation.

The key insight is that the intermediate reasoning steps, while beneficial
for learning, need not be produced as discrete tokens at inference time.
Instead, we train the model through a curriculum that gradually replaces
textual reasoning steps with special \emph{latent tokens} whose embeddings
are populated by the model's own hidden states, effectively internalising
the reasoning process.

Our approach comprises three components:
(1)~vocabulary extension and embedding initialisation
    (\S\ref{sec:vocab});
(2)~a multi-pass forward mechanism with hidden-state replacement that
    enables reasoning in continuous space (\S\ref{sec:coconut_forward});
    and
(3)~a stage-based curriculum that orchestrates the transition from
    explicit to latent reasoning (\S\ref{sec:curriculum}).

\subsection{Vocabulary Extension and Embedding Initialisation}
\label{sec:vocab}

Starting from the pre-trained Qwen2.5-3B-Instruct model, we augment its
tokeniser with three groups of task-specific tokens:
(i)~514 \emph{motion tokens} (\texttt{<Motion>}, \texttt{</Motion>}, and
\texttt{<Motion\_0>} through \texttt{<Motion\_511>}) that establish a
one-to-one mapping with the VQ-VAE codebook;
(ii)~2 \emph{think markers} (\texttt{<think>}, \texttt{</think>}) that
delimit explicit reasoning regions; and
(iii)~3 \emph{latent tokens} (\texttt{<|start-latent|>},
\texttt{<|end-latent|>}, \texttt{<|latent|>}) that demarcate the
continuous reasoning region.
The model's embedding matrix and language modelling head are jointly
resized to accommodate the expanded vocabulary.
Following the initialisation strategy of
Coconut~\cite{hao2024training}, we initialise the \texttt{<|latent|>}
embedding from that of ``\texttt{<<}''---a token whose semantic role as a
delimiter provides a more informative starting point than random
initialisation.

\subsection{Multi-Pass Forward with Latent Token Replacement}
\label{sec:coconut_forward}

The central mechanism of our approach replaces the standard single-pass
forward computation with a multi-pass procedure that iteratively fills
latent token embeddings with contextualised hidden states.
Concretely, let
$\mathbf{x}{=}(\mathbf{q},\,t_{\mathrm{s}},\,
  \underbrace{l_1,\ldots,l_M}_{\text{latent}},\,
  t_{\mathrm{e}},\,\mathbf{r},\,\mathbf{a})$
denote an input sequence containing $M$ latent tokens, where $\mathbf{q}$
is the question, $t_{\mathrm{s}}/t_{\mathrm{e}}$ are the latent-region
boundary markers, $\mathbf{r}$ represents any remaining explicit reasoning
steps, and $\mathbf{a}$ is the answer.
The forward computation proceeds as follows.

\paragraph{Iterative hidden-state injection (passes $1,\ldots,M$).}
At the $i$-th pass, the model performs a forward computation from the end
of the previous pass up to the position of the $i$-th latent token $l_i$,
reusing the KV cache accumulated in prior passes to avoid redundant
computation.
The last-layer hidden state $\mathbf{h}_{i-1}$ at the position
immediately preceding $l_i$ is then extracted and used to overwrite the
embedding of $l_i$:
%
\begin{equation}
  \mathbf{e}(l_i) \;\leftarrow\; \mathbf{h}_{i-1}.
  \label{eq:latent_replace}
\end{equation}
%
This operation effectively \emph{conditions} each latent token on the
model's own contextualised representation of all preceding tokens,
allowing information to propagate through the latent region via successive
hidden-state injections rather than through discrete token generation.

\paragraph{Final pass.}
After all $M$ latent embeddings have been replaced, a final forward pass
processes the remaining tokens $(\mathbf{r}, \mathbf{a})$ using the full
accumulated KV cache, producing logits over the vocabulary at each
position.

\paragraph{Loss computation.}
The logits from all $M{+}1$ passes are concatenated along the sequence
dimension:
%
\begin{equation}
  \hat{\mathbf{y}}
  = \big[\hat{\mathbf{y}}^{(1)};\;\hat{\mathbf{y}}^{(2)};\;\ldots\;;\;
    \hat{\mathbf{y}}^{(M+1)}\big],
  \label{eq:logit_concat}
\end{equation}
%
and the model is optimised with the standard next-token cross-entropy
objective:
%
\begin{equation}
  \mathcal{L}
  = -\!\sum_{j\in\mathcal{A}}
    \log\, p_\theta\!\left(x_{j+1}\mid\hat{\mathbf{y}}_j\right),
  \label{eq:loss}
\end{equation}
%
where $\mathcal{A}$ indexes the set of \emph{unmasked} positions.
The question tokens, latent tokens (including boundary markers), and any
padding are assigned label $-100$ and excluded from the loss, so that the
model is supervised exclusively on the remaining reasoning tokens (when
present) and the answer motion tokens.

\paragraph{Efficient batched computation.}
Two design choices ensure computational efficiency.
First, at each pass $i>1$, the KV cache from the preceding pass is sliced
to the appropriate prefix length and reused, avoiding recomputation of
attention over already-processed positions.
Second, we employ a \emph{left-padding collation} strategy that aligns the
first \texttt{<|latent|>} position across all samples in a mini-batch,
enabling the $M$ replacement passes to operate at uniform positions and
thus permitting fully batched multi-pass forward without padding-induced
waste.

\subsection{Curriculum Learning Schedule}
\label{sec:curriculum}

Training the model to reason entirely in latent space from scratch is
intractable, as the latent tokens carry no semantic content at
initialisation.
We therefore adopt a stage-based curriculum that begins with fully explicit
reasoning and gradually increases the proportion of latent tokens.

Formally, let $E_s$ denote the number of training epochs per curriculum
stage and $K_{\max}$ the maximum number of reasoning steps eligible for
replacement.
At epoch $e$, the scheduled curriculum stage is:
%
\begin{equation}
  k(e) = \min\!\Big(\big\lfloor e\,/\,E_s \big\rfloor,\;K_{\max}\Big).
  \label{eq:stage}
\end{equation}
%
At stage $k$, the first $k$ reasoning steps $(s_1,\ldots,s_k)$ are
excised from the input and replaced by $k{\times}c$ latent tokens, where
$c$ denotes the number of latent tokens allocated per replaced step.
The remaining steps $(s_{k+1},\ldots,s_N)$ are retained as explicit text:
%
\begin{equation}
  \mathbf{x}^{(k)}
  = \big(\mathbf{q},\;t_{\mathrm{s}},\;
    \underbrace{l_1,\ldots,l_{kc}}_{\text{latent}},\;
    t_{\mathrm{e}},\;
    \underbrace{s_{k+1}\cdots s_N}_{\text{explicit CoT}},\;
    \mathbf{a}\big).
  \label{eq:curriculum_seq}
\end{equation}
%
Once $k$ reaches $K_{\max}$, all explicit reasoning is removed and the
latent region is zero-padded to a fixed length of $K_{\max}{\cdot}c$
tokens, yielding a fully latent input.
This gradual transition provides the model with a structured learning
signal: in early stages the explicit reasoning tokens act as
``scaffolding'' that guides the latent representations, while in later
stages the model must rely entirely on its internalised continuous
reasoning.

\paragraph{Stage randomisation for regularisation.}
To prevent the model from overfitting to the rigid stage boundaries and to
encourage robustness across varying ratios of latent-to-explicit tokens,
we introduce stochastic stage sampling: with probability $p_u$ each
training sample independently draws a random stage
$k'{\sim}\mathrm{Uniform}(0, N_i)$ in lieu of the scheduled stage $k(e)$,
where $N_i$ is the number of available reasoning steps for sample $i$.
This acts as a form of data augmentation in the curriculum dimension.

Table~\ref{tab:curriculum} summarises the curriculum progression under our
default hyperparameter setting.

\begin{table}[t]
\centering
\caption{Curriculum schedule with $E_s{=}3$, $K_{\max}{=}8$, $c{=}2$.
  The model transitions from fully explicit CoT to fully latent reasoning
  over 30 training epochs.}
\label{tab:curriculum}
\setlength{\tabcolsep}{6pt}
\begin{tabular}{cccp{4.6cm}}
\toprule
Epochs & Stage $k$ & Latent tokens & Supervised target \\
\midrule
0--2   & 0 & 0  & Full CoT + motion tokens \\
3--5   & 1 & 2  & Steps $(s_2,\ldots)$ + motion \\
6--8   & 2 & 4  & Steps $(s_3,\ldots)$ + motion \\
$\vdots$ & $\vdots$ & $\vdots$ & $\vdots$ \\
21--23 & 7 & 14 & Steps $(s_8,\ldots)$ + motion \\
24--29 & ${\geq}8$ & 16 & Motion tokens only \\
\bottomrule
\end{tabular}
\end{table}

\subsection{Inference}
\label{sec:generation}

At test time, the input is constructed with the maximum number of latent
tokens ($K_{\max}{\cdot}c{=}16$) and no explicit reasoning steps,
corresponding to the fully latent regime of the final curriculum stage.
Generation proceeds in two phases:

\noindent\textbf{Phase~1: Latent reasoning.}
The multi-pass forward mechanism (Sec.~\ref{sec:coconut_forward})
processes all $M$ latent tokens, iteratively injecting hidden-state
representations via Eq.~\eqref{eq:latent_replace}.
Upon completion, the model's KV cache encodes a continuous ``chain of
thought'' that captures the reasoning necessary for motion generation,
without producing any intermediate text tokens.

\noindent\textbf{Phase~2: Autoregressive motion decoding.}
Conditioned on the KV cache from Phase~1, the model autoregressively
generates motion tokens via greedy decoding until either an end-of-sequence
token is emitted or the maximum generation budget (256 tokens) is
exhausted.

Compared to explicit CoT inference, this two-phase procedure replaces $T$
sequential text-generation steps with $M$ incremental KV-cached forward
passes, where typically $M{\ll}T$.
We provide a detailed computational analysis of this trade-off in
Sec.~\ref{sec:cost_analysis}.

\subsection{Computational Cost Analysis}
\label{sec:cost_analysis}

We now provide a formal comparison of the inference-time computational
cost between explicit CoT and latent CoT reasoning.
Consider a Transformer decoder with $L$ layers, hidden dimension~$d$,
and attention head dimension~$d_k$ with $n_h$ heads.
Let $|\mathbf{q}|$ denote the question length, $T$ the number of
explicit thinking tokens in a text CoT chain, $M{=}K_{\max}{\cdot}c$ the
number of latent tokens, and $M_a$ the number of motion tokens to
generate.

\paragraph{Unit cost model.}
Under KV-cache-enabled autoregressive decoding, the cost of one
incremental forward step (processing a single new token against a
cache of length~$S$) is dominated by the feed-forward computation for
moderate sequence lengths ($S\ll d^2/d_k$):
%
\begin{equation}
  F_{\mathrm{step}}
  = 2L\bigl(C_{\mathrm{attn}}(S) + C_{\mathrm{ffn}}\bigr)
  \approx 24\,L\,d^2,
  \label{eq:step_cost}
\end{equation}
%
where $C_{\mathrm{attn}}(S){=}4d^2 + 2n_h d_k S$ accounts for
QKV projections and the attention dot product against the cached keys,
and $C_{\mathrm{ffn}}{\approx}8d^2$ for a standard two-layer FFN\@.
A full forward pass over $n$ tokens without KV cache costs approximately
$n\,F_{\mathrm{step}}$, since each position incurs comparable FFN cost
and the quadratic attention term $O(n^2 d_k n_h)$ remains subdominant
for the sequence lengths typical in our setting
($n{\leq}512,\;d{=}3072$).

\paragraph{Explicit CoT inference.}
The model first prefills the question ($|\mathbf{q}|$ tokens), then
autoregressively generates $T$ thinking tokens followed by $M_a$ motion
tokens:
%
\begin{equation}
  C_{\mathrm{explicit}}
  = \underbrace{|\mathbf{q}|\,F_{\mathrm{step}}}_{\text{prefill}}
  \;+\;
  \underbrace{(T + M_a)\,F_{\mathrm{step}}}_{\text{sequential decoding}}.
  \label{eq:cost_explicit}
\end{equation}

\paragraph{Latent CoT inference.}
With KV cache reuse across the $M$ iterative hidden-state replacement
passes, the latent reasoning phase reduces to $M$ incremental forward
steps (each processing one latent position against the accumulated
cache) rather than $M$ full-sequence passes.
The total cost comprises three terms:
%
\begin{equation}
  C_{\mathrm{latent}}
  = \underbrace{|\mathbf{q}|\,F_{\mathrm{step}}}_{\text{prefill}}
  \;+\;
  \underbrace{(M + 1)\,F_{\mathrm{step}}}_{\text{latent passes}}
  \;+\;
  \underbrace{M_a\,F_{\mathrm{step}}}_{\text{motion decoding}},
  \label{eq:cost_latent}
\end{equation}
%
where the $+1$ accounts for the end-of-latent boundary marker.

\paragraph{Reasoning-phase speedup.}
Since the motion decoding cost $M_a\,F_{\mathrm{step}}$ is shared by both
approaches, the inference advantage of latent CoT derives entirely from
the reasoning phase.
The speedup ratio over explicit CoT reasoning is:
%
\begin{equation}
  \rho = \frac{T}{M + 1}.
  \label{eq:speedup}
\end{equation}
%
In our setting, $M{=}K_{\max}{\cdot}c{=}16$ latent tokens replace an
average of $T{\approx}63$ text thinking tokens
(${\sim}4.2$~steps~$\times$~${\sim}15$~tokens/step), yielding
$\rho{\approx}3.7{\times}$.

\paragraph{Discussion.}
The speedup is governed by the \emph{compression ratio} $T/M$: how many
text tokens each latent token effectively replaces.
This compression is possible because latent representations reside in
continuous $\mathbb{R}^d$ space ($d{=}3072$) rather than in a discrete
vocabulary of size $|\mathcal{V}|$, enabling each hidden state to encode
richer intermediate computations than a single text token.
As noted by Hao~et~al.~\cite{hao2024training}, continuous reasoning can
maintain multiple hypotheses simultaneously (akin to breadth-first
search), whereas text CoT is constrained to a single serial reasoning
chain (depth-first search).

However, for tasks where the explicit reasoning chain is short
($T{\lesssim}M$), the multi-pass overhead approaches or exceeds that of
direct text generation, and the latent formulation confers no speed
advantage.
More broadly, the primary motivation for latent CoT in our setting is
not inference speed per se, but rather the capacity of continuous
representations to capture motion-relevant reasoning structure that may
be difficult to express in discrete token form.

\subsection{Implementation Details}
\label{sec:training_details}

We fine-tune Qwen2.5-3B-Instruct in BFloat16 precision using the AdamW
optimiser.
Training is distributed across multiple GPUs via PyTorch
DistributedDataParallel (DDP); we found that FSDP is incompatible with our
multi-pass forward mechanism due to the variable number of forward calls
per sample triggering NCCL synchronisation deadlocks.
To stabilise training under mixed precision, we apply gradient norm
clipping at 1.0\@.
The optimiser state is reinitialised at each curriculum stage transition to
prevent stale momentum from the previous stage's different loss landscape.

Table~\ref{tab:hyperparams} provides the complete hyperparameter
configuration.

\begin{table}[t]
\centering
\caption{Training hyperparameters.}
\label{tab:hyperparams}
\setlength{\tabcolsep}{10pt}
\begin{tabular}{ll}
\toprule
Hyperparameter & Value \\
\midrule
Base model & Qwen2.5-3B-Instruct \\
Precision & BFloat16 \\
Parallelism & DDP \\
Optimiser & AdamW \\
Learning rate & $1\times10^{-4}$ \\
Weight decay & $1\times10^{-2}$ \\
LR warm-up & 200 steps (linear) \\
Gradient clipping & max norm 1.0 \\
Per-GPU batch size & 16 \\
Gradient accumulation & 8 steps \\
Max sequence length & 512 \\
Total epochs & 30 \\
Epochs per stage ($E_s$) & 3 \\
Max latent stage ($K_{\max}$) & 8 \\
Latent tokens per step ($c$) & 2 \\
Stage randomisation ($p_u$) & 0.1 \\
Motion codebook size ($K$) & 512 \\
\bottomrule
\end{tabular}
\end{table}

\paragraph{Validation protocol.}
We monitor two complementary evaluation signals during training.
\emph{Validation loss} is computed every epoch on a held-out subset of
100~samples using the curriculum dataset of the current stage.
\emph{Generation evaluation} is conducted every 3~epochs: the model
generates motion token sequences in the fully latent regime, which are
then compared against ground-truth sequences via two token-level metrics:
(1)~\emph{exact match}---the fraction of samples for which the predicted
and reference motion token sequences are identical; and
(2)~\emph{token accuracy}---the proportion of correctly predicted tokens
in the overlapping prefix, normalised by the ground-truth sequence length.


% ============================================================================
% BibTeX entries (add to your .bib file)
% ============================================================================
% @article{hao2024training,
%   title     = {Training Large Language Models to Reason in a Continuous
%                Latent Space},
%   author    = {Hao, Shibo and Gu, Sainbayar and Ma, Haotian and
%                Hong, Joshua Jiahua and Wang, Zhen and Wang, Daisy Zhe
%                and Hu, Zhiting},
%   journal   = {arXiv preprint arXiv:2412.06769},
%   year      = {2024}
% }
%
% @inproceedings{guo2022generating,
%   title     = {Generating Diverse and Natural {3D} Human Motions from Text},
%   author    = {Guo, Chuan and Zou, Shihao and Zuo, Xinxin and Wang, Sen
%                and Ji, Tao and Li, Xingyu and Cheng, Li},
%   booktitle = {CVPR},
%   year      = {2022}
% }
%
% @inproceedings{zhang2023generating,
%   title     = {{T2M-GPT}: Generating Human Motion from Textual
%                Descriptions with Discrete Representations},
%   author    = {Zhang, Jianrong and Zhang, Yangsong and Cun, Xiaodong
%                and Huang, Shaoli and Zhang, Yong and Zhao, Hongwei
%                and Lu, Hongtao and Shen, Xi},
%   booktitle = {CVPR},
%   year      = {2023}
% }
